% ============================================================
% Chapter 1: Introduction
% ============================================================

\chapter{Introduction}

\section{Introduction}

Personal Finance Manager (PFM) is an AI-powered hybrid mobile application designed to simplify expense tracking through natural language processing. In an era where financial literacy and expense management are crucial for personal well-being, traditional expense tracking applications often fail to gain user adoption due to their cumbersome interfaces requiring manual data entry through complex forms and rigid categorization systems.

PFM addresses these challenges by allowing users to record expenses using natural language input. Instead of filling out multiple form fields, users can simply type ``biryani rahul 500'' or ``taxi to airport 1500,'' and the system automatically extracts the amount, categorizes the expense, and identifies any associated parties. This approach significantly reduces the friction of expense tracking, encouraging consistent usage and better financial awareness.

The application leverages cutting-edge technologies including Google Gemini AI for natural language understanding, React for the user interface, FastAPI for the backend, and Supabase for database and authentication services. The mobile experience is delivered through Capacitor, enabling native Android deployment from a single codebase.

\section{Problem Statement}

Managing personal finances effectively requires consistent tracking of income and expenses. However, existing solutions present several challenges:

\begin{enumerate}
    \item \textbf{Tedious Data Entry}: Traditional expense tracking applications require users to fill out multiple form fields (amount, category, description, date) for each transaction, discouraging regular use.
    
    \item \textbf{Manual Categorization}: Users must manually select appropriate categories for each expense, which is time-consuming and prone to inconsistency.
    
    \item \textbf{Lack of Intelligence}: Most applications do not leverage AI to understand context or provide intelligent insights about spending patterns.
    
    \item \textbf{Limited Query Capabilities}: Users cannot easily query their expense data using natural language questions like ``How much did I spend on food this month?''
    
    \item \textbf{Poor Loan Tracking}: Tracking money lent to or borrowed from friends and family is often overlooked in personal finance applications.
\end{enumerate}

\section{Objectives}

The primary objectives of this project are:

\begin{enumerate}
    \item \textbf{Natural Language Expense Entry}: Develop a system that accepts expense input in natural language and automatically extracts relevant information (amount, item, category, person).
    
    \item \textbf{Intelligent Categorization}: Implement automatic expense categorization using a hybrid approach combining regex pattern matching and AI-powered parsing.
    
    \item \textbf{Multi-Group Support}: Enable users to manage expenses across different groups (personal, family, shared expenses).
    
    \item \textbf{Comprehensive Loan Tracking}: Provide functionality to track money lent and borrowed, with automatic categorization of loan-related transactions.
    
    \item \textbf{RAG-Based Querying}: Implement a Retrieval-Augmented Generation system allowing users to query their expenses using natural language.
    
    \item \textbf{Visual Analytics}: Provide interactive charts and insights for financial analysis and spending pattern visualization.
    
    \item \textbf{Cross-Platform Deployment}: Develop a responsive web application that can be deployed as a native Android application.
\end{enumerate}

\section{Scope and Limitation}

\subsection{Scope}

The scope of this project includes:

\begin{itemize}
    \item Natural language processing for expense parsing
    \item Automatic expense categorization (20+ categories)
    \item User authentication and authorization
    \item Multi-group expense management
    \item Loan tracking (money lent and borrowed)
    \item Income tracking
    \item RAG-based natural language querying
    \item Visual analytics with interactive charts
    \item Web and Android deployment
    \item Real-time data synchronization
\end{itemize}

\subsection{Limitations}

The current implementation has the following limitations:

\begin{itemize}
    \item \textbf{Language Support}: Currently optimized for English with some support for Nepali currency formats
    \item \textbf{Offline Mode}: Requires internet connectivity for AI-powered parsing and data synchronization
    \item \textbf{Bank Integration}: Does not integrate with banking APIs for automatic transaction import
    \item \textbf{Multi-Currency}: Limited support for multiple currencies
    \item \textbf{iOS Support}: Currently deployed only on Android; iOS deployment is planned for future releases
\end{itemize}

\section{Development Methodology}

This project follows an \textbf{Agile/Iterative Development} methodology, which emphasizes:

\begin{itemize}
    \item \textbf{Incremental Development}: Features are developed and tested incrementally, allowing for early feedback and continuous improvement.
    
    \item \textbf{Iterative Refinement}: The NLP parsing system was iteratively refined based on testing with various input formats.
    
    \item \textbf{Test-Driven Approach}: Extensive testing scripts were developed to validate NLP parsing accuracy.
    
    \item \textbf{Continuous Integration}: Regular integration of frontend and backend components with automated testing.
\end{itemize}

The development process consisted of the following phases:

\begin{enumerate}
    \item \textbf{Requirements Analysis}: Understanding user needs and defining functional requirements
    \item \textbf{System Design}: Designing the architecture, database schema, and API endpoints
    \item \textbf{Implementation}: Developing frontend and backend components
    \item \textbf{Testing}: Unit testing, integration testing, and user acceptance testing
    \item \textbf{Deployment}: Deploying the application on web and Android platforms
\end{enumerate}

\section{Report Organization}

This report is organized into the following chapters:

\begin{itemize}
    \item \textbf{Chapter 1: Introduction} --- Provides an overview of the project, problem statement, objectives, scope, and development methodology.
    
    \item \textbf{Chapter 2: Background Study and Literature Review} --- Discusses the theoretical foundations, related technologies, and review of similar systems.
    
    \item \textbf{Chapter 3: System Analysis and Design} --- Presents the requirement analysis, feasibility study, and system design including UML diagrams.
    
    \item \textbf{Chapter 4: Implementation and Testing} --- Details the implementation of various modules, tools used, and testing procedures.
    
    \item \textbf{Chapter 5: Conclusion and Future Recommendations} --- Summarizes the project outcomes and suggests future enhancements.
\end{itemize}
